\documentclass[../../../include/open-logic-section]{subfiles}

\begin{document}

\olfileid[fr]{sth}{story}{predicative}	

\olsection{Prédicatif et imprédicatif}

L'ensemble de Russell, $R$, était défini par $\Setabs{x}{x \notin
x}$. Plus explicitement, $R$ serait l'ensemble qui contiendrait tous
les ensembles qui ne s'appartiennent pas à eux-mêmes, et eux seuls.
En définissant $R$, nous quantifions donc sur un domaine qui
contiendrait $R$ (s'il existait).

Il s'agit d'une définition \emph{imprédicative}. Plus généralement,
on peut dire qu'une définition est imprédicative si et seulement si
elle quantifie sur un domaine qui contient l'objet défini.

À la suite des paradoxes, Whitehead, Russell, Poincar\'{e} et Weyl
ont rejeté ces définitions imprédicatives comme relevant d'un «~cercle
vicieux~» :
\begin{quote}
	Une analyse des paradoxes à éviter montre qu'ils résultent tous
	d'une forme de cercle vicieux. Les cercles vicieux en question
	naissent de la supposition qu'une collection d'objets puisse
	contenir des membres qui ne peuvent être définis qu'au moyen de
	la collection prise comme un tout [\ldots. \textparagraph]

	Le principe qui nous permet d'éviter les totalités illégitimes
	peut s'énoncer ainsi : «~Ce qui fait intervenir la \emph{totalité}
	d'une collection ne doit pas être un membre de cette collection~» ;
	ou, réciproquement : «~Si une certaine collection, à supposer
	qu'elle forme une totalité, devait comporter des membres qui ne
	puissent être définis qu'au moyen de cette totalité, alors cette
	collection ne forme pas une totalité.~» Nous appellerons ce
	principe le «~principe du cercle vicieux~», parce qu'il permet
	d'éviter les cercles vicieux qu'entraîne la supposition de
	totalités illégitimes.
	\citep[p.~37]{WhiteheadRussell1910}
\end{quote}
Si nous suivons ces auteurs en adoptant le \emph{principe du cercle
vicieux}, nous pourrions essayer de remplacer le désastreux schéma de
compréhension naïve (de la \olref[sth][story][rus]{sec}) par un
principe de ce genre :

\begin{defish}
\emph{Compréhension prédicative.} Pour toute formule $\phi$ qui ne quantifie que sur les ensembles, l'ensemble$^\prime$ $\Setabs{x}{\phi(x)}$ existe.
\end{defish}

Tant que les ensembles$^{\prime}$ ne sont pas des ensembles, aucune
contradiction ne s'ensuit.

Malheureusement, la compréhension prédicative n'est guère
\emph{exhaustive}. Après tout, elle introduit de nouvelles entités :
les ensembles$^{\prime}$. Il faut donc aussi considérer les formules
qui quantifient sur les ensembles$^{\prime}$. Si ces formules donnent
toujours un ensemble$^{\prime}$, le paradoxe de Russell réapparaît :
il suffit de considérer l'ensemble$^{\prime}$ de tous les
ensembles$^{\prime}$ qui ne s'appartiennent pas à eux-mêmes. Dans la
même logique, il faut dire qu'une formule qui quantifie sur les
ensembles$^{\prime}$ donne un ensemble$^{\prime\prime}$ correspondant.
Il nous faudra ensuite des ensembles$^{\prime\prime\prime}$, des
ensembles$^{\prime\prime\prime\prime}$, et ainsi de suite. Pour éviter
une prolifération de primes, il sera plus simple de parler
d'ensembles$_0$, d'ensembles$_1$, d'ensembles$_2$, d'ensembles$_3$,
d'ensembles$_4$,\ldots. Cela nous conduit vers la théorie (simple)
des types.

Une telle théorie soulève quelques objections assez évidentes (même
s'il n'est pas évident qu'elles soient \emph{décisives}). En bref : la
théorie obtenue est lourde à manier ; elle multiplie les sortes
d'objets qu'elle postule ; et, finalement, il n'est pas certain que
les définitions imprédicatives soient \emph{si problématiques que
cela}.

Pour éclairer ce dernier point, considérons cette remarque de
\citeauthor{Ramsey1925} :
\begin{quote}
	nous pouvons désigner un homme comme le plus grand d'un groupe, et
	l'identifier ainsi au moyen d'une totalité dont il est lui-même
	membre, sans qu'il y ait pour autant de cercle vicieux.
	\citep{Ramsey1925}
\end{quote}
Pour Ramsey, «~l'homme le plus grand du groupe~» \emph{est} une
définition imprédicative ; elle est pourtant manifestement
parfaitement acceptable.

On pourrait répondre que, dans ce cas, nous pourrions désigner la
personne la plus grande par des moyens \emph{prédicatifs}. Nous
pourrions par exemple simplement la montrer du doigt. L'objection
contre les définitions imprédicatives devrait alors clairement être
limitée aux entités qui ne peuvent être désignées \emph{que} de manière
imprédicative. Mais même dans ce cas, il faudrait expliquer davantage
pourquoi une telle «~imprédicativité essentielle~» serait si
problématique.\footnote{Pour aller plus loin, voir \citet{Linnebo2010}.}

Certes, les définitions imprédicatives posent un problème majeur si
nous attendons d'une définition qu'elle nous fournisse une sorte de
recette pour \emph{créer} un objet. Une définition imprédicative nous
enfermerait alors réellement dans un cercle vicieux : pour créer
l'objet ainsi défini, il faudrait \emph{d'abord} créer tous les objets
(y compris celui-là), puisque sa définition fait intervenir leur
totalité ; il faudrait donc créer cet objet avant même de l'avoir
créé. Mais, là encore, il ne s'agit d'une objection sérieuse contre
les ensembles définis de façon «~essentiellement imprédicative~» que
si nous considérons les ensembles comme des choses que nous
\emph{créons}. Et ce n'est (probablement) pas notre point de vue.

Ainsi, pour le meilleur ou pour le pire, l'approche devenue courante
ne consiste pas à adopter une position stricte sur la
(im)\-pré\-di\-ca\-ti\-vi\-té. Elle repose plutôt sur ce que l'on
appelle aujourd'hui l'approche cumulative et itérative. Celle-ci nous
permettra finalement de répartir les ensembles en «~étapes~», un
\emph{peu} comme l'approche prédicative répartit les entités entre
ensembles$_0$, ensembles$_1$, ensembles$_2$, \ldots, mais sans
postuler de différence de nature entre eux.
\footnote{Note de traduction : la source décrit par erreur les ensembles qui ne sont pas non auto-appartenants, alors que la formule exige $x\notin x$ ; elle dit aussi de rejeter le principe du cercle vicieux, alors que le passage cité et le schéma suivant l'adoptent. Ces deux inversions sont corrigées ici.}
\end{document}
